\section{Observable transfer theorem for the physical $n=2$ mode}
\label{sec:bhsm-observable-transfer}

The reduced BHSM mode becomes observationally useful only after its
action-generated phase-space data are mapped to metric and matter response
variables.  This section defines that map without introducing an additional
degree of freedom.

Let
\begin{equation}
X_2(t)
=
\begin{pmatrix}
q_2(t)\\
\Pi_2(t)
\end{pmatrix}
\end{equation}
denote the physical $n=2$ phase-space state.  Linear response of the retained
metric/matter variables can be represented abstractly as
\begin{equation}
Y_2(t)
=
\mathbb T_2(t,t_i)X_2(t_i)
+
\int_{t_i}^{t}
\mathbb G_2(t,t')\,S_2(t')\,dt',
\label{eq:bhsm-linear-transfer}
\end{equation}
where $Y_2$ may contain the two scalar metric potentials, matter density and
velocity response, and any auxiliary response coordinate used by the chosen
four-dimensional representation.  The operators $\mathbb T_2$ and
$\mathbb G_2$ are determined by the reduced action and its source projection;
they are not independent phenomenological fields.

\subsection{Angular-factorization theorem}

\paragraph{Theorem 2 (single-mode observable factorization).}
Assume (i) the retained perturbation is the pure $A_{4i}$ component of the
physical $n=2$ scalar-topographic mode, (ii) the response is linear, and
(iii) the transfer operator is rotationally covariant on the closed
background.  Then every scalar linear observable $\mathcal O$ sourced by this
single mode has the leading angular form
\begin{equation}
\boxed{
\delta\mathcal O(z,\hat{\mathbf n})
=
\mathcal A_{\mathcal O}(z)
\left(
\hat{\mathbf n}\!\cdot\!\hat{\mathbf p}
\right)
},
\label{eq:observable-factorization}
\end{equation}
where the entire observable dependence is contained in the redshift transfer
amplitude $\mathcal A_{\mathcal O}(z)$ and the sky axis
$\hat{\mathbf p}$ is shared by all observables at leading one-mode order.

\paragraph{Proof.}
The physical perturbation belongs to one $n=2$ representation of the
isometry group of $S^3$.  A rotationally covariant linear transfer operator
cannot generate a different harmonic representation in the absence of an
additional anisotropic source.  Restriction of the retained $A_{4i}$
subspace to the observer shell is exactly
$T_{2,\rm dip}\propto\hat{\mathbf n}\cdot\hat{\mathbf p}$ by
Eq.~\eqref{eq:bhsm-n2-exact-dipole}.  Applying any scalar linear observation
functional changes only its radial/time coefficient.  Therefore the angular
factor is common and Eq.~\eqref{eq:observable-factorization} follows.
\hfill$\square$

\subsection{Observable kernels}

For each observable define a linear functional
$\mathcal L_{\mathcal O,z}$ acting on Eq.~\eqref{eq:bhsm-linear-transfer}.
Then
\begin{equation}
\mathcal A_{\mathcal O}(z)
=
\mathcal L_{\mathcal O,z}
\left[
\mathbb T_2X_2(t_i)
+
\int\mathbb G_2S_2\,dt'
\right].
\label{eq:observable-amplitude-functional}
\end{equation}

The principal late-time observables used in this work are represented by
\begin{align}
\frac{\delta D_L}{D_L}
&=
\mathcal A_{D_L}(z)\,
(\hat{\mathbf n}\!\cdot\!\hat{\mathbf p}),
\\
\frac{\delta H}{H}
&=
\mathcal A_H(z)\,
(\hat{\mathbf n}\!\cdot\!\hat{\mathbf p}),
\\
\frac{\delta D_A}{D_A}
&=
\mathcal A_{D_A}(z)\,
(\hat{\mathbf n}\!\cdot\!\hat{\mathbf p}),
\\
\delta(\Phi+\Psi)
&=
\mathcal A_{\rm lens}(z)\,
(\hat{\mathbf n}\!\cdot\!\hat{\mathbf p}),
\\
\frac{\delta(f\sigma_8)}{f\sigma_8}
&=
\mathcal A_{f\sigma_8}(z)\,
(\hat{\mathbf n}\!\cdot\!\hat{\mathbf p}),
\end{align}
whenever the corresponding observable is evaluated in the linear response
regime.

For supernova distance modulus,
\begin{equation}
\delta\mu
=
\frac{5}{\ln 10}
\frac{\delta D_L}{D_L},
\end{equation}
so that
\begin{equation}
\boxed{
A_\mu(z)
=
\frac{5}{\ln 10}\,
\mathcal A_{D_L}(z)
}.
\label{eq:sn-transfer}
\end{equation}
Equation~\eqref{eq:sn-transfer} supplies the direct bridge between the
geometric $n=2$ carrier and the low-redshift supernova dipole amplitude.

\subsection{Shared-axis and amplitude-ratio predictions}

Once the transfer kernels are fixed by the reduced action, the amplitudes are
not independent fit parameters.  For any two linear observables,
\begin{equation}
\frac{
\mathcal A_{\mathcal O_1}(z)
}{
\mathcal A_{\mathcal O_2}(z)
}
=
\frac{
\mathcal L_{\mathcal O_1,z}[\mathcal R_2]
}{
\mathcal L_{\mathcal O_2,z}[\mathcal R_2]
},
\end{equation}
where $\mathcal R_2$ denotes the common physical response generated by
$X_2$ and its projected sources.  Thus the single-mode branch predicts a
common leading axis and linked redshift-dependent amplitude ratios across
distance, expansion, growth, and lensing observables.

A statistically persistent axis mismatch between independent observables,
after survey-window and calibration effects are controlled, would falsify
the pure one-mode realization.  Likewise, a leading angular morphology that
cannot be represented by the $A_{4i}$ dipole subspace requires additional
physical harmonics or rejection of the one-mode branch.

\subsection{Higher-mode control}

For the harmonic filter
\begin{equation}
C_n
=
c_s^2-\chi\left[1-\ell^2n(n+2)\right],
\end{equation}
the interval
\begin{equation}
\frac1{15}<\ell^2<\frac18
\end{equation}
softens the physical $n=2$ carrier while stiffening every $n\ge3$ scalar
harmonic.  The next source-normalized response obeys
\begin{equation}
\left|
\frac{T_3/S_3}{T_2/S_2}
\right|
=
\frac{8C_2}{15C_3}.
\end{equation}
This provides the explicit control parameter for departures from the
single-axis theorem.

\subsection{Closure criterion for the manuscript}

Scientific closure of the observable bridge requires the reduced operator
$\mathcal K_2$ to determine the transfer kernels entering
Eq.~\eqref{eq:observable-amplitude-functional}.  A four-dimensional EFT
realization is acceptable only if it reproduces the same two-dimensional
reduced phase space $(q_2,\Pi_2)$ and the same observable transfer kernels
without adding an independent physical scalar.  Reference-branch numerical
curves may be quoted before this exact identification only when they are
explicitly labeled as reference predictions rather than framework-wide
theorems.
